\scanpage{149}
% lines 1–5
Posons $\mu_\alpha=\mu'_\alpha+\mu''_\alpha$

$\mu'_\alpha$ partie multiplicative $=$ nb de branches \uline{formelles} passant par $P_\alpha$, \deleted{\unsure{de}} \unsure{diminué} de 1

$\mu''_\alpha$ partie additive \deleted{\unsure{de}}.

\[
\dim G_{\mathrm{mult}}=\sum\mu'_\alpha-c_0+1=h_{\underline{1}}
\]
% slanted
\margin{\uline{invariant} \uline{topologique} \uline{de la} \uline{situation}}

% lines 6–15
[N.B. Lorsque les $C_i$ \deleted{\unsure{forment|passent} \unsure{deux à deux} \unsure{en des} pts multiples \unsure{ordinaires}}, alors $h_1$ \deleted{\unsure{se}} est bien déterminé \deleted{par} combinatoirement par le schéma \deleted{\unsure{de rencontre}} \inserted{combinatoire de $C$}. Il est \unsure{en ts cas} déterminé par le schéma combinatoire de la \unsure{donnée|situation} $\widetilde C\to C$ ; c'est aussi le rang du \deleted{\unsure{\(\pi_1\)}} noyau de $\deleted{\unsure{\pi_1}}H^1(C)\to\deleted{\unsure{\prod\pi_1|\prod}}H^1(\widetilde C)$, (Cohomologie de Weil) \unsure{ou} du conoyau de
\[
\prod\pi_1(\widetilde C_i)\unsure{-|\to}\pi_1(C)\ \big]
\]
\[
\dim G_{\mathrm{add}}=\sum\mu''_\alpha
\]

% lines 16–19
Dire que $\mathrm{Pic}^0$ n'a que des composantes \deleted{\unsure{c'est|$C_i$}} multiplicatives, signifie que les \unsure{recollements|multiplicités} en $P_\alpha$ sont triviaux \inserted{formellement; Alors $\mu_\alpha=\mu'_\alpha\geq2$,}. Dire que c'est une V.A. \deleted{signifie} \inserted{implique} que \deleted{\unsure{les}} de plus
% archive sheet number: 148

